% !TEX program  = lualatex
% Cheat Sheet — alle Umgebungen & Befehle der Lernzettel-Vorlage
% Kompilieren: latexmk -lualatex cheatsheet.tex  (oder: make cheatsheet)
% ──────────────────────────────────────────────────────────────────────────────
\documentclass[a4paper, 11pt, ngerman]{article}

\usepackage[table, dvipsnames]{xcolor}

%% colormeta.tex — Alle Design-Tokens: Farben, Schriften, Abstände
%% ══════════════════════════════════════════════════════════════════════════
%% NUR DIESE DATEI bearbeiten, um das gesamte Dokument umzugestalten.
%% Keine Farben oder Schriften sind anderswo hart kodiert.
%% ══════════════════════════════════════════════════════════════════════════

%% ── Schriften (werden von fontspec in styles.tex gesetzt) ─────────────────
%% Standard: TeX Gyre-Schriften (immer in TeX Live enthalten, keine
%% Installation erforderlich). Für schönere Alternativen siehe README.md.
\newcommand{\mainfont}{Arial}                % Standard-Schriftart
\newcommand{\sansfont}{Arial}                % Standard-Schriftart
\newcommand{\monofont}{TeX Gyre Cursor}     % Courier-Klon, funktional

%% Optionale Premium-Schriften (Homebrew/manuell installieren):
%% \newcommand{\mainfont}{Source Serif 4}
%% \newcommand{\sansfont}{Source Sans 3}
%% \newcommand{\monofont}{JetBrains Mono}

%% ── Abstände ──────────────────────────────────────────────────────────────
\newcommand{\contentmargin}{2.5cm}   % Seitenrand ringsum
\newcommand{\parskipval}{6pt}        % Abstand zwischen Absätzen
\newcommand{\boxinnersep}{10pt}      % Innenabstand von Info-/Warnboxen

%% ── Farben ────────────────────────────────────────────────────────────────
%% Voraussetzung: xcolor ist bereits mit [table,dvipsnames] in main.tex geladen.

%% Brand / Thema
\definecolor{primary}{HTML}{2E4057}         % Dunkelblau  — Überschriften, Tabellenköpfe
\definecolor{accent}{HTML}{048A81}          % Petrol/Teal — Links, Akzente

%% Infobox (blau-türkis)
\definecolor{infocolor}{HTML}{DDF3F5}       % Hintergrund
\definecolor{infoframe}{HTML}{048A81}       % Rahmen (= accent)

%% Warnbox (amber)
\definecolor{warncolor}{HTML}{FFF4E0}       % Hintergrund
\definecolor{warnframe}{HTML}{E07B39}       % Rahmen

%% Pro / Contra
\definecolor{procolor}{HTML}{D6EDDA}        % Grün-Hintergrund
\definecolor{proframe}{HTML}{1E8A3F}        % Grün-Rahmen
\definecolor{concolor}{HTML}{FADDE1}        % Rot-Hintergrund
\definecolor{conframe}{HTML}{C0392B}        % Rot-Rahmen

%% Codeblock
\definecolor{codeblock}{HTML}{F5F5F5}       % Hintergrund
\definecolor{codeframe}{HTML}{CCCCCC}       % Rahmen

%% Tabellen
\definecolor{tableheader}{HTML}{2E4057}     % Kopfzeilenhintergrund (= primary)
\definecolor{tableheadertext}{HTML}{FFFFFF} % Kopfzeilentext
\definecolor{tablerowalt}{HTML}{EEF2F7}     % Alternierend eingefärbte Zeilen

%% Prozess-Schritte
\definecolor{stepcolor}{HTML}{048A81}       % Schrittnummer-Badge (= accent)

%% Merksatz (Indigo/Lila)
\definecolor{merksatzcolor}{HTML}{F0EEFF}   % Hintergrund
\definecolor{merksatzframe}{HTML}{5C4EC2}   % Rahmen

%% Beispiel (Smaragdgrün)
\definecolor{beispielcolor}{HTML}{E8F5E9}   % Hintergrund
\definecolor{beispielframe}{HTML}{2E9E68}   % Rahmen

%% Aufgabe (Goldgelb)
\definecolor{aufgabecolor}{HTML}{FFF8E1}    % Hintergrund
\definecolor{aufgabeframe}{HTML}{E6A817}    % Rahmen

%% Konzeptkarte
\definecolor{konzeptcolor}{HTML}{F7F9FC}    % Hintergrund (blau-grau)

%% usermeta.tex — Persönliche Metadaten
%% Hier Autor, Kurs, Matrikelnummer, Semester und Datum eintragen.
%% ──────────────────────────────────────────────────────────────────────────

\newcommand{\docauthor}{Levin Rüßmann}
\newcommand{\doccourse}{Informatik B.Sc.}
\newcommand{\docstudentid}{838791}
\newcommand{\docsemester}{Sommersemester 2026}
\newcommand{\docdate}{\today}

%% topicmeta.tex — Dokumenten- / Themenmetadaten
%% Hier Titel, Untertitel, Fach und Beschreibung eintragen.
%% ──────────────────────────────────────────────────────────────────────────

\newcommand{\doctitle}{Cloudtechnologie }
\newcommand{\docsubtitle}{Zusammenfassung und Lernzettel}
\newcommand{\docsubject}{Grundlagen der Cloudtechnologie}
\newcommand{\docdescription}{%
  Kompakte Zusammenfassung der Vorlesungsinhalte für die Prüfungsvorbereitung.%
}

%% packages.tex — Alle \usepackage-Deklarationen
%% Reihenfolge beachten: xcolor ist bereits in main.tex geladen.
%% Konfiguration (die colormeta-Makros verwendet) erfolgt in styles.tex.
%% ──────────────────────────────────────────────────────────────────────────

%% ── Schriften & Kodierung (LuaLaTeX) ──────────────────────────────────────
\usepackage{fontspec}          % Systemschriften über OpenType/TrueType
\usepackage{unicode-math}      % Unicode-Mathematikschriften

%% ── Sprache ───────────────────────────────────────────────────────────────
\usepackage[ngerman]{babel}
\usepackage{csquotes}          % kontextsensitive Anführungszeichen (für biblatex)

%% ── Seitenlayout ──────────────────────────────────────────────────────────
\usepackage{geometry}          % Seitenränder; konfiguriert in styles.tex
\usepackage{parskip}           % Absatzabstand statt Einzug
\usepackage{setspace}

%% ── Typographie ───────────────────────────────────────────────────────────
\usepackage{microtype}         % optischer Randausgleich, Unterschneidung

%% ── Mathematik ────────────────────────────────────────────────────────────
\usepackage{amsmath}
\usepackage{amssymb}

%% ── Grafiken & Gleitumgebungen ────────────────────────────────────────────
\usepackage{graphicx}
\usepackage{float}
\usepackage{caption}
\usepackage{subcaption}

%% ── Tabellen ──────────────────────────────────────────────────────────────
\usepackage{booktabs}          % professionelle Linien (\toprule, \midrule…)
\usepackage{tabularx}          % auto-breite Spalten (X-Typ)
\usepackage{array}             % erweiterte Spaltentypen
\usepackage{multirow}          % Zellen über mehrere Zeilen
\usepackage{colortbl}          % farbige Tabellenzeilen/-zellen

%% ── Listen ────────────────────────────────────────────────────────────────
\usepackage{enumitem}          % flexible Listen-Konfiguration

%% ── Boxen & Rahmen ────────────────────────────────────────────────────────
\usepackage[most]{tcolorbox}   % bunte Boxen (infobox, warnbox, process…)

%% ── Listings / Code ───────────────────────────────────────────────────────
\usepackage{listings}

%% ── Kopf- & Fußzeile ──────────────────────────────────────────────────────
\usepackage{fancyhdr}

%% ── Überschriften & TOC ───────────────────────────────────────────────────
\usepackage{titlesec}          % Überschriftenformatierung
\usepackage[titles]{tocloft}   % TOC-Formatierung

%% ── Werkzeuge für eigene Befehle ──────────────────────────────────────────
\usepackage{etoolbox}          % \appto, \preto usw. (für procon-Umgebung)
\usepackage{xparse}            % \NewDocumentEnvironment (ggf. im Kernel)
\usepackage{environ}           % \NewEnviron: Body-Kollektion für tabularx-Wrapper

%% ── Bibliographie ─────────────────────────────────────────────────────────
\usepackage[
  backend    = biber,
  style      = authoryear,
  sortlocale = de_DE,
  maxnames   = 3,
  minnames   = 1,
]{biblatex}
\addbibresource{bibliography/sources.bib}

%% ── Hyperlinks (muss fast zuletzt geladen werden) ─────────────────────────
\usepackage{hyperref}          % konfiguriert in styles.tex via \hypersetup
\usepackage{bookmark}

%% ── Diverses ──────────────────────────────────────────────────────────────
\usepackage{lipsum}            % Blindtext für Beispiele

%% commands.tex — Eigene Umgebungen und Befehle
%% Alle Farben kommen aus config/colormeta.tex; nichts ist hart kodiert.
%% ──────────────────────────────────────────────────────────────────────────

%% ══════════════════════════════════════════════════════════════════════════
%% 1.  process  +  \step{text}
%%     Nummerierter Schritt-für-Schritt-Ablauf mit farbigem linken Rand.
%%
%%  Verwendung:
%%    \begin{process}
%%      \step{Erster Schritt}
%%      \step{Zweiter Schritt}
%%    \end{process}
%% ══════════════════════════════════════════════════════════════════════════
\newcounter{procstep}

\newenvironment{process}{%
  \setcounter{procstep}{0}%
  \begin{tcolorbox}[
    enhanced,
    breakable,
    boxrule         = 0pt,
    frame hidden,
    borderline west = {3pt}{0pt}{stepcolor},
    colback         = stepcolor!6!white,
    left            = 10pt,
    right           = 8pt,
    top             = 6pt,
    bottom          = 6pt,
    arc             = 3pt,
  ]%
}{%
  \end{tcolorbox}%
}

\newcommand{\step}[1]{%
  \stepcounter{procstep}%
  \vspace{4pt}%
  \noindent
  \colorbox{stepcolor}{%
    \color{white}\bfseries\small\hspace{3pt}\theprocstep\hspace{3pt}%
  }%
  \hspace{6pt}#1\par
}

%% ══════════════════════════════════════════════════════════════════════════
%% 2.  infobox[Titel]
%%     Farbige Infobox mit optionalem fettem Titel.
%%
%%  Verwendung:
%%    \begin{infobox}[Mein Titel] ... \end{infobox}
%%    \begin{infobox}             ... \end{infobox}   % ohne Titel
%% ══════════════════════════════════════════════════════════════════════════
\NewDocumentEnvironment{infobox}{O{}}{%
  \def\lzinfoboxtitle{#1}%
  \ifdefempty{\lzinfoboxtitle}{%
    %% ── Ohne Titel ──
    \begin{tcolorbox}[
      enhanced, breakable,
      colback  = infocolor,
      colframe = infoframe,
      boxrule  = 1pt,
      arc      = 4pt,
      left     = \boxinnersep,
      right    = \boxinnersep,
      top      = \boxinnersep,
      bottom   = \boxinnersep,
    ]%
  }{%
    %% ── Mit Titel ──
    \begin{tcolorbox}[
      enhanced, breakable,
      colback      = infocolor,
      colframe     = infoframe,
      boxrule      = 1pt,
      arc          = 4pt,
      left         = \boxinnersep,
      right        = \boxinnersep,
      top          = \boxinnersep,
      bottom       = \boxinnersep,
      title        = {\textbf{\lzinfoboxtitle}},
      fonttitle    = \bfseries\color{white},
      colbacktitle = infoframe,
      attach boxed title to top left = {yshift=-2mm, xshift=4mm},
      boxed title style = {
        colback  = infoframe,
        colframe = infoframe,
        arc = 3pt, top = 2pt, bottom = 2pt,
      },
    ]%
  }%
}{%
  \end{tcolorbox}%
}

%% ══════════════════════════════════════════════════════════════════════════
%% 3.  warnbox[Titel]
%%     Warnbox / Tipp-Box mit optionalem fettem Titel.
%% ══════════════════════════════════════════════════════════════════════════
\NewDocumentEnvironment{warnbox}{O{}}{%
  \def\lzwarnboxtitle{#1}%
  \ifdefempty{\lzwarnboxtitle}{%
    %% ── Ohne Titel ──
    \begin{tcolorbox}[
      enhanced, breakable,
      colback  = warncolor,
      colframe = warnframe,
      boxrule  = 1pt,
      arc      = 4pt,
      left     = \boxinnersep,
      right    = \boxinnersep,
      top      = \boxinnersep,
      bottom   = \boxinnersep,
    ]%
  }{%
    %% ── Mit Titel ──
    \begin{tcolorbox}[
      enhanced, breakable,
      colback      = warncolor,
      colframe     = warnframe,
      boxrule      = 1pt,
      arc          = 4pt,
      left         = \boxinnersep,
      right        = \boxinnersep,
      top          = \boxinnersep,
      bottom       = \boxinnersep,
      title        = {\textbf{\lzwarnboxtitle}},
      fonttitle    = \bfseries\color{white},
      colbacktitle = warnframe,
      attach boxed title to top left = {yshift=-2mm, xshift=4mm},
      boxed title style = {
        colback  = warnframe,
        colframe = warnframe,
        arc = 3pt, top = 2pt, bottom = 2pt,
      },
    ]%
  }%
}{%
  \end{tcolorbox}%
}

%% ══════════════════════════════════════════════════════════════════════════
%% 4.  procon  +  \pro{text}  \con{text}
%%     Zwei-Spalten-Gegenüberstellung Vorteile / Nachteile.
%%     Alle \pro{} und \con{} werden gesammelt und am Ende nebeneinander
%%     als zwei tcolorbox-Minipages gesetzt.
%%
%%  Verwendung:
%%    \begin{procon}
%%      \pro{Vorteil 1}
%%      \con{Nachteil 1}
%%    \end{procon}
%% ══════════════════════════════════════════════════════════════════════════
\newenvironment{procon}{%
  \def\lzproitems{}%
  \def\lzconitems{}%
  % \def statt \newcommand: kein Fehler beim mehrfachen Aufruf
  % ##1 = #1 auf zweiter Makro-Ebene innerhalb von \newenvironment
  \def\pro##1{\appto\lzproitems{\item ##1}}%
  \def\con##1{\appto\lzconitems{\item ##1}}%
}{%
  \par\smallskip\noindent
  \begin{minipage}[t]{0.475\linewidth}
    \begin{tcolorbox}[
      enhanced,
      colback      = procolor,
      colframe     = proframe,
      boxrule      = 1pt,
      arc          = 4pt,
      title        = {Vorteile},
      fonttitle    = \bfseries\color{white},
      colbacktitle = proframe,
      left = 6pt, right = 6pt, top = 4pt, bottom = 4pt,
    ]
      \begin{itemize}[leftmargin=*, nosep, topsep=0pt]
        \lzproitems
      \end{itemize}
    \end{tcolorbox}
  \end{minipage}%
  \hfill
  \begin{minipage}[t]{0.475\linewidth}
    \begin{tcolorbox}[
      enhanced,
      colback      = concolor,
      colframe     = conframe,
      boxrule      = 1pt,
      arc          = 4pt,
      title        = {Nachteile},
      fonttitle    = \bfseries\color{white},
      colbacktitle = conframe,
      left = 6pt, right = 6pt, top = 4pt, bottom = 4pt,
    ]
      \begin{itemize}[leftmargin=*, nosep, topsep=0pt]
        \lzconitems
      \end{itemize}
    \end{tcolorbox}
  \end{minipage}%
  \par\smallskip
}

%% ══════════════════════════════════════════════════════════════════════════
%% 5.  codebox
%%     Monospace-Block mit grauem Hintergrund (für kurze Ausschnitte).
%%     Für Syntax-Highlighting: listings-Paket nutzen (siehe packages.tex).
%% ══════════════════════════════════════════════════════════════════════════
\newenvironment{codebox}{%
  \begin{tcolorbox}[
    enhanced,
    breakable,
    colback  = codeblock,
    colframe = codeframe,
    boxrule  = 0.5pt,
    arc      = 3pt,
    left     = 8pt,
    right    = 8pt,
    top      = 6pt,
    bottom   = 6pt,
    fontupper = \ttfamily\small,
  ]%
}{%
  \end{tcolorbox}%
}

%% ══════════════════════════════════════════════════════════════════════════
%% 6.  \sidetables{Inhalt links}{Inhalt rechts}
%%     Zwei beliebige Inhalte (z. B. Tabellen, Listen) nebeneinander.
%% ══════════════════════════════════════════════════════════════════════════
\newcommand{\sidetables}[2]{%
  \par\smallskip\noindent
  \begin{minipage}[t]{0.475\linewidth}
    #1
  \end{minipage}%
  \hfill
  \begin{minipage}[t]{0.475\linewidth}
    #2
  \end{minipage}%
  \par\smallskip
}

%% ══════════════════════════════════════════════════════════════════════════
%% 7.  stdtable{Spaltendefinition}{Kopfzeile}
%%     Formatierte Tabelle mit farbiger Kopfzeile und alternierenden Zeilen.
%%     Spaltentypen: l r c  für Text; X für auto-breite Spalten; p{…} für
%%     feste Breite. X-Spalten sind durch tabularx möglich, weil \NewEnviron
%%     (environ-Paket) den Body komplett sammelt, bevor tabularx ihn scannt.
%%
%%  Verwendung:
%%    \begin{stdtable}{l X X}{Algorithmus & Kompl. & Notiz}
%%      Bubble Sort & $O(n^2)$ & Einfach \\
%%    \end{stdtable}
%% ══════════════════════════════════════════════════════════════════════════
\NewEnviron{stdtable}[2]{%
  \par\smallskip
  \rowcolors{2}{tablerowalt}{white}%
  \begin{center}%
  \begin{tabularx}{\linewidth}{#1}%
  \rowcolor{tableheader}%
  \color{tableheadertext}\bfseries #2\\[\jot]%
  \BODY
  \end{tabularx}%
  \end{center}%
  \rowcolors{1}{}{}%   Reset: verhindert Übertrag auf nachfolgende Tabellen
  \par\smallskip
}

%% styles.tex — Schriften, Seitenlayout, Kopf-/Fußzeile, TOC, Hyperlinks
%% Alle Farben und Schriften kommen aus config/colormeta.tex.
%% ──────────────────────────────────────────────────────────────────────────

%% ── Schriften (fontspec) ──────────────────────────────────────────────────
\setmainfont{\mainfont}[
  Ligatures = TeX,
]
\setsansfont{\sansfont}[
  Ligatures = TeX,
]
\setmonofont{\monofont}[
  Scale = 0.88,
]

%% ── Seitenlayout (geometry) ───────────────────────────────────────────────
\geometry{
  a4paper,
  top    = \contentmargin,
  bottom = \contentmargin,
  left   = \contentmargin,
  right  = \contentmargin,
  headheight = 15pt,
}

%% ── Absatzformat ─────────────────────────────────────────────────────────
\setlength{\parskip}{\parskipval}
\setlength{\parindent}{0pt}

%% ── Grafikpfad ───────────────────────────────────────────────────────────
\graphicspath{{assets/images/}}

%% ── Überschriften (titlesec) ─────────────────────────────────────────────
\AddToHook{cmd/section/before}{\clearpage}

\titleformat{\section}
  {\Large\bfseries\color{primary}}         % Format
  {\color{primary}\thesection}             % Label
  {0.75em}                                  % Abstand Label–Text
  {}                                        % Vor dem Titel
  [\color{primary!50!white}\rule{\linewidth}{0.6pt}]  % Unter dem Titel

\titleformat{\subsection}
  {\large\bfseries\color{primary!80!black}}
  {\color{accent}\thesubsection}
  {0.75em}
  {}

\titleformat{\subsubsection}
  {\normalsize\bfseries\color{accent}}
  {\thesubsubsection}
  {0.75em}
  {}

\titlespacing*{\section}      {0pt}{1.4ex plus .4ex minus .2ex}{0.6ex}
\titlespacing*{\subsection}   {0pt}{1.2ex plus .3ex minus .2ex}{0.4ex}
\titlespacing*{\subsubsection}{0pt}{1.0ex plus .2ex minus .1ex}{0.3ex}

%% ── Inhaltsverzeichnis (tocloft) ─────────────────────────────────────────
\renewcommand{\cftsecfont}{\bfseries\color{primary}}
\renewcommand{\cftsecpagefont}{\bfseries\color{primary}}
\renewcommand{\cftsubsecfont}{\color{primary!80!black}}
\renewcommand{\cftsubsecpagefont}{\color{primary!80!black}}
\setlength{\cftsecindent}{0pt}
\setlength{\cftsubsecindent}{1.5em}
\setlength{\cftbeforesecskip}{2pt}

%% ── Kopf- und Fußzeile (fancyhdr) ────────────────────────────────────────
\pagestyle{fancy}
\fancyhf{}
\fancyhead[L]{\small\color{primary}\docsubject}
\fancyhead[R]{\small\color{primary}\docauthor}
\fancyfoot[C]{\small\thepage}
\renewcommand{\headrulewidth}{0.4pt}
\renewcommand{\footrulewidth}{0pt}
% Auch Kapitelanfangsseiten mit fancyhdr versehen
\fancypagestyle{plain}{%
  \fancyhf{}%
  \fancyfoot[C]{\small\thepage}%
  \renewcommand{\headrulewidth}{0pt}%
}

%% ── Hyperlinks (hyperref) ────────────────────────────────────────────────
\hypersetup{
  colorlinks = true,
  linkcolor  = primary,
  urlcolor   = accent,
  citecolor  = accent,
  pdftitle   = {\doctitle},
  pdfauthor  = {\docauthor},
  pdfsubject = {\docsubject},
}

%% ── Captions ─────────────────────────────────────────────────────────────
\captionsetup{
  font       = small,
  labelfont  = {bf, color=primary},
  labelsep   = period,
}

%% ── Titelseite (wird in main.tex mit \makelernzetteltitle aufgerufen) ─────
\newcommand{\makelernzetteltitle}{%
  \begin{titlepage}
    \centering
    \vspace*{2.5cm}

    % Fach-Badge
    {\sffamily\small\color{white}%
      \colorbox{accent}{\hspace{6pt}\docsubject\hspace{6pt}}%
    }\par
    \vspace{1.2cm}

    % Haupttitel
    {\Huge\bfseries\color{primary}\doctitle\par}
    \vspace{0.5cm}
    {\large\color{accent}\docsubtitle\par}

    \vspace{0.8cm}
    % Trennlinie
    {\color{primary!40!white}\rule{0.6\linewidth}{1.5pt}}\par
    \vspace{0.5cm}

    % Beschreibung
    {\normalsize\color{primary!70!black}\docdescription\par}

    \vfill

    % Metadaten-Tabelle
    \begin{tabular}{@{}r@{\hspace{8pt}}l@{}}
      \color{primary!60!black}\small Autor            & \docauthor      \\[2pt]
      \color{primary!60!black}\small Matrikelnummer   & \docstudentid   \\[2pt]
      \color{primary!60!black}\small Studiengang      & \doccourse      \\[2pt]
      \color{primary!60!black}\small Semester         & \docsemester    \\[2pt]
      \color{primary!60!black}\small Datum            & \docdate        \\
    \end{tabular}

    \vspace{2cm}
    % Fußleiste
    {\color{primary!30!white}\rule{\linewidth}{0.6pt}}\par
    \vspace{0.3cm}
    {\small\color{primary!50!black}\sffamily Lernzettel}
  \end{titlepage}%
  \clearpage
}


% Engere Ränder als im Hauptdokument
\geometry{a4paper, margin=1.4cm, top=1.6cm, bottom=1.6cm,
          includeheadfoot, headheight=14pt}

% LaTeX-Syntax-Highlighting für tcblisting
\lstdefinestyle{latexcs}{%
  language=[LaTeX]TeX,
  basicstyle=\ttfamily\scriptsize,
  keywordstyle=\color{accent}\bfseries,
  commentstyle=\color{primary!45!white}\itshape,
  texcsstyle=*\color{accent},
  moretexcs={step,pro,con,sidetables,makelernzetteltitle,
             docauthor,doccourse,docstudentid,docsemester,doctitle,
             docsubtitle,docsubject,docdescription,mainfont,sansfont,
             monofont,contentmargin,parskipval,boxinnersep,definecolor,
             rowcolor,rowcolors,cellcolor},
  breaklines=true,
  frame=none,
  keepspaces=true,
  showstringspaces=false,
  columns=flexible,
  tabsize=2,
}

% Basis-Style für alle Cheatsheet-Boxen
\tcbset{
  cs/.style={
    enhanced, breakable,
    colback=white, colframe=primary!18!white,
    boxrule=0.7pt, arc=5pt,
    fonttitle=\sffamily\bfseries\small,
    colbacktitle=primary, coltitle=white,
    top=6pt, bottom=6pt, left=7pt, right=7pt,
    listing options={style=latexcs},
  },
}

% Kopf- & Fußzeile des Cheatsheets
\pagestyle{fancy}
\fancyhf{}
\fancyhead[L]{\small\sffamily\color{primary}\textbf{Lernzettel-Template} \textcolor{primary!45!black}{— Cheat Sheet}}
\fancyhead[R]{\small\sffamily\color{primary!55!black}\today}
\fancyfoot[C]{\small\color{primary!50!black}\thepage}
\renewcommand{\headrulewidth}{0.5pt}
\renewcommand{\headrule}{\color{primary!25!white}\hrule width\headwidth height\headrulewidth}

% ─────────────────────────────────────────────────────────────────────────────
\begin{document}
\thispagestyle{empty}

% ── Großer Header ─────────────────────────────────────────────────────────────
\begin{tcolorbox}[
  enhanced, arc=6pt, boxrule=0pt,
  colback=primary, coltext=white,
  left=16pt, right=16pt, top=13pt, bottom=13pt,
]
  {\huge\bfseries\sffamily Lernzettel-Template}%
  \quad{\normalsize\color{accent!70!white}\sffamily v1.0 · LuaLaTeX + Biber}\\[0.3em]
  {\large\color{accent!80!white} Cheat Sheet — Alle Umgebungen \& Befehle auf einen Blick}\\[0.25em]
  {\small\color{white!65!primary}
    Konfiguration: \texttt{config/colormeta.tex}\enspace|\enspace
    Kompilieren: \texttt{make build}\enspace|\enspace
    Dokumentation: \texttt{README.md}
  }
\end{tcolorbox}

\vspace{0.55em}

% ── 1 + 2: infobox / warnbox nebeneinander ───────────────────────────────────
\noindent
\begin{minipage}[t]{0.487\linewidth}
\begin{tcblisting}{cs,
  title={\texttt{infobox} — Informations-/Definitions-Box},
  listing above text,
}
% Ohne Titel:
\begin{infobox}
  Allgemeiner Hinweis oder Info.
\end{infobox}
% Mit optionalem Titel:
\begin{infobox}[Definition: Algorithmus]
  Eine präzise, endliche Folge von
  Schritten zur Lösung eines Problems.
\end{infobox}
\end{tcblisting}
\end{minipage}%
\hfill
\begin{minipage}[t]{0.487\linewidth}
\begin{tcblisting}{cs,
  title={\texttt{warnbox} — Warnhinweis / Tipp-Box},
  listing above text,
}
% Ohne Titel:
\begin{warnbox}
  Wichtiger Hinweis ohne Titel.
\end{warnbox}
% Mit optionalem Titel:
\begin{warnbox}[Achtung: Off-by-one!]
  Array-Indizes beginnen bei 0,
  nicht bei 1. Typische Fehlerquelle.
\end{warnbox}
\end{tcblisting}
\end{minipage}

\vspace{0.45em}

% ── 3: process + \step ────────────────────────────────────────────────────────
\begin{tcblisting}{cs,
  title={\texttt{process} + \texttt{\textbackslash step\{text\}} — Nummerierte Schrittfolge mit farbigem Rand},
  listing side text,
  righthand width=0.41\linewidth,
}
\begin{process}
  \step{Aufgabe lesen und verstehen}
  \step{Lösungsstrategie entwerfen}
  \step{Implementieren und testen}
  \step{Ergebnisse dokumentieren}
\end{process}
\end{tcblisting}

\vspace{0.45em}

% ── 4: codebox ────────────────────────────────────────────────────────────────
\begin{tcblisting}{cs,
  title={\texttt{codebox} — Monospace-Block, grauer Hintergrund (kein Syntax-Highlighting)},
  listing side text,
  righthand width=0.41\linewidth,
}
\begin{codebox}
latexmk -lualatex main.tex\\
biber main\\
latexmk -lualatex main.tex
\end{codebox}
\end{tcblisting}

\vspace{0.45em}

% ── 5: procon ─────────────────────────────────────────────────────────────────
\begin{tcblisting}{cs,
  title={\texttt{procon} + \texttt{\textbackslash pro\{...\}} / \texttt{\textbackslash con\{...\}} — Vor-/Nachteil-Vergleich, zwei Spalten},
  listing above text,
}
\begin{procon}
  \pro{Klare visuelle Trennung von Vorteilen und Nachteilen}
  \pro{Beliebig viele Einträge pro Seite möglich}
  \con{Nur für Stichpunkt-Vergleiche, kein Fließtext}
  \con{Breite fest bei je 47{,}5\,\% der Textbreite}
\end{procon}
\end{tcblisting}

\vspace{0.45em}

% ── 6: stdtable ───────────────────────────────────────────────────────────────
\begin{tcblisting}{cs,
  title={\texttt{stdtable\{Spalten\}\{Kopfzeile\}} — Formatierte Tabelle, farbige Kopfzeile, alternierende Zeilen},
  listing above text,
}
\begin{stdtable}{l X r}{Algorithmus & Beschreibung & Komplexität}
  Quicksort  & Teile-und-Herrsche, in-place          & $O(n \log n)$ \\
  Bubblesort & Vergleichsbasiert, einfach             & $O(n^2)$      \\
  Mergesort  & Stabil, rekursiv, extra Speicher       & $O(n \log n)$ \\
\end{stdtable}
\end{tcblisting}

\vspace{0.45em}

% ── 7: \sidetables ────────────────────────────────────────────────────────────
\begin{tcblisting}{cs,
  title={\texttt{\textbackslash sidetables\{Links\}\{Rechts\}} — Zwei beliebige Inhalte nebeneinander (Tabellen, Listen, …)},
  listing above text,
}
\sidetables{%
  \begin{stdtable}{l X}{Sprache & Paradigma}
    Python & Interpretiert, multi-paradigm \\
    Rust   & Kompiliert, speichersicher    \\
  \end{stdtable}
}{%
  \begin{stdtable}{l X}{Datenstruktur & Zugriff}
    Array  & $O(1)$ Random-Access \\
    Liste  & $O(n)$ lineare Suche \\
  \end{stdtable}
}
\end{tcblisting}

% ── 8 + 9 + 10: merksatz / beispiel / aufgabe ────────────────────────────────
\noindent
\begin{minipage}[t]{0.318\linewidth}
\begin{tcblisting}{cs,
  title={\texttt{merksatz} — Wichtiger Satz},
  listing above text,
}
% Mit Titel:
\begin{merksatz}[Satz]
  $a^2 + b^2 = c^2$
\end{merksatz}
% Ohne Titel:
\begin{merksatz}
  Immer terminieren!
\end{merksatz}
\end{tcblisting}
\end{minipage}%
\hfill
\begin{minipage}[t]{0.318\linewidth}
\begin{tcblisting}{cs,
  title={\texttt{beispiel} — Beispiel-Box},
  listing above text,
}
\begin{beispiel}[Quicksort]
  Pivot = 3:\\
  Links: [1, 1]\\
  Rechts: [4, 5]
\end{beispiel}
\end{tcblisting}
\end{minipage}%
\hfill
\begin{minipage}[t]{0.318\linewidth}
\begin{tcblisting}{cs,
  title={\texttt{aufgabe} — Aufgabe (gezählt)},
  listing above text,
}
\begin{aufgabe}[Analyse]
  Zeitkomplexität von
  Bubblesort bestimmen.
\end{aufgabe}
\begin{aufgabe}
  Zähler läuft automatisch.
\end{aufgabe}
\end{tcblisting}
\end{minipage}

\vspace{0.45em}

% ── 11: konzeptkarte + \kzweispaltig ─────────────────────────────────────────
\begin{tcblisting}{cs,
  title={\texttt{konzeptkarte\{Titel\}} + \texttt{\textbackslash kzweispaltig\{links\}\{rechts\}} — Karte mit Tabelle \& Prozess nebeneinander},
  listing above text,
}
\begin{konzeptkarte}{Mergesort}
  Stabiler Teile-und-Herrsche-Algorithmus. Garantiert $O(n \log n)$.
  \kzweispaltig{%
    \begin{stdtable}{l X}{Eigenschaft & Wert}
      Stabil  & Ja  \\
      Speicher & $O(n)$ \\
    \end{stdtable}
  }{%
    \begin{process}
      \step{Liste halbieren}
      \step{Rekursiv sortieren}
      \step{Zusammenmischen}
    \end{process}
  }
\end{konzeptkarte}
\end{tcblisting}

\vspace{0.45em}

% ── 12 + 13: \icode + syntaxbox ──────────────────────────────────────────────
\noindent
\begin{minipage}[t]{0.36\linewidth}
\begin{tcblisting}{cs,
  title={\texttt{\textbackslash icode\{...\}} — Inline-Code},
  listing above text,
}
Die Funktion \icode{sort(arr)}
hat Komplexität $O(n \log n)$.
\icode{NullPointerException}
tritt bei \icode{arr == null}
auf.
\end{tcblisting}
\end{minipage}%
\hfill
\begin{minipage}[t]{0.61\linewidth}
\begin{tcblisting}{cs,
  title={\texttt{syntaxbox[Sprache]} — Code mit Syntax-Highlighting (Standard: Python)},
  listing above text,
}
\begin{syntaxbox}[Python]
def quicksort(arr):
    if len(arr) <= 1:
        return arr
    pivot = arr[len(arr) // 2]
    left  = [x for x in arr if x < pivot]
    right = [x for x in arr if x > pivot]
    return quicksort(left) + [pivot] + quicksort(right)
\end{syntaxbox}
\end{tcblisting}
\end{minipage}

\clearpage

% ── 14: Konfiguration ─────────────────────────────────────────────────────────
\section*{Konfiguration}

\noindent
\begin{minipage}[t]{0.487\linewidth}

\begin{tcblisting}{cs,
  title={\texttt{config/usermeta.tex} — Autordaten},
  listing only,
}
% Diese Werte erscheinen auf der Titelseite.
\newcommand{\docauthor}{Max Mustermann}
\newcommand{\doccourse}{Informatik B.Sc.}
\newcommand{\docstudentid}{12345678}
\newcommand{\docsemester}{SoSe 2025}
\newcommand{\docdate}{\today}
\end{tcblisting}

\vspace{0.4em}

\begin{tcblisting}{cs,
  title={\texttt{config/topicmeta.tex} — Thema \& Titel},
  listing only,
}
\newcommand{\doctitle}{Grundlagen der Algorithmen}
\newcommand{\docsubtitle}{Zusammenfassung und Lernzettel}
\newcommand{\docsubject}{Algorithmen und Datenstrukturen}
\newcommand{\docdescription}{%
  Prüfungsvorbereitung für Semester 3.%
}
\end{tcblisting}

\end{minipage}%
\hfill
\begin{minipage}[t]{0.487\linewidth}

\begin{tcblisting}{cs,
  title={\texttt{config/colormeta.tex} — Alle Design-Tokens (Auszug)},
  listing only,
}
%% Schriften (kein Install nötig, in TeX Live enthalten)
\newcommand{\mainfont}{TeX Gyre Pagella}
\newcommand{\sansfont}{TeX Gyre Heros}
\newcommand{\monofont}{TeX Gyre Cursor}

%% Abstände
\newcommand{\contentmargin}{2.5cm}
\newcommand{\parskipval}{6pt}
\newcommand{\boxinnersep}{10pt}

%% Farben — NUR hier anpassen, nirgendwo sonst!
\definecolor{primary}{HTML}{2E4057}    % Dunkelblau
\definecolor{accent}{HTML}{048A81}     % Petrol/Teal
\definecolor{infocolor}{HTML}{DDF3F5}
\definecolor{infoframe}{HTML}{048A81}
\definecolor{warncolor}{HTML}{FFF4E0}
\definecolor{warnframe}{HTML}{E07B39}
% ... weitere Farben (procolor, concolor, etc.)
\end{tcblisting}

\end{minipage}

\vspace{0.55em}

% ── 9: Farbpalette ────────────────────────────────────────────────────────────
\section*{Farbpalette (\texttt{colormeta.tex})}

\begin{tcolorbox}[cs, title={Alle definierten Farben — nur \texttt{colormeta.tex} anpassen}]
\renewcommand{\arraystretch}{2.2}
\setlength{\tabcolsep}{4pt}
\begin{tabular}{@{} p{3.55cm} p{3.55cm} p{3.55cm} p{3.55cm} @{}}
  \cellcolor{primary}\textcolor{white}{\texttt{\textbf{primary}}}%
    \newline\textcolor{white}{\scriptsize\texttt{\#2E4057 · Dunkelblau}} &
  \cellcolor{accent}\textcolor{white}{\texttt{\textbf{accent}}}%
    \newline\textcolor{white}{\scriptsize\texttt{\#048A81 · Petrol/Teal}} &
  \cellcolor{infocolor}\textcolor{primary}{\texttt{\textbf{infocolor}}}%
    \newline\textcolor{primary}{\scriptsize\texttt{\#DDF3F5 · Info-BG}} &
  \cellcolor{infoframe}\textcolor{white}{\texttt{\textbf{infoframe}}}%
    \newline\textcolor{white}{\scriptsize\texttt{\#048A81 · Info-Rand}}
  \\
  \cellcolor{warncolor}\textcolor{primary}{\texttt{\textbf{warncolor}}}%
    \newline\textcolor{primary}{\scriptsize\texttt{\#FFF4E0 · Warn-BG}} &
  \cellcolor{warnframe}\textcolor{white}{\texttt{\textbf{warnframe}}}%
    \newline\textcolor{white}{\scriptsize\texttt{\#E07B39 · Warn-Rand}} &
  \cellcolor{procolor}\textcolor{primary}{\texttt{\textbf{procolor}}}%
    \newline\textcolor{primary}{\scriptsize\texttt{\#D6EDDA · Pro-BG}} &
  \cellcolor{proframe}\textcolor{white}{\texttt{\textbf{proframe}}}%
    \newline\textcolor{white}{\scriptsize\texttt{\#1E8A3F · Pro-Rand}}
  \\
  \cellcolor{concolor}\textcolor{primary}{\texttt{\textbf{concolor}}}%
    \newline\textcolor{primary}{\scriptsize\texttt{\#FADDE1 · Con-BG}} &
  \cellcolor{conframe}\textcolor{white}{\texttt{\textbf{conframe}}}%
    \newline\textcolor{white}{\scriptsize\texttt{\#C0392B · Con-Rand}} &
  \cellcolor{codeblock}\textcolor{primary}{\texttt{\textbf{codeblock}}}%
    \newline\textcolor{primary}{\scriptsize\texttt{\#F5F5F5 · Code-BG}} &
  \cellcolor{codeframe}\textcolor{primary}{\texttt{\textbf{codeframe}}}%
    \newline\textcolor{primary}{\scriptsize\texttt{\#CCCCCC · Code-Rand}}
  \\
  \cellcolor{tableheader}\textcolor{white}{\texttt{\textbf{tableheader}}}%
    \newline\textcolor{white}{\scriptsize\texttt{\#2E4057 · Tabelle-Kopf}} &
  \cellcolor{tablerowalt}\textcolor{primary}{\texttt{\textbf{tablerowalt}}}%
    \newline\textcolor{primary}{\scriptsize\texttt{\#EEF2F7 · Tabelle-Zeile}} &
  \cellcolor{stepcolor}\textcolor{white}{\texttt{\textbf{stepcolor}}}%
    \newline\textcolor{white}{\scriptsize\texttt{\#048A81 · Schritt-Badge}} &
  \cellcolor{white}\textcolor{primary!40!white}{\textit{\small weitere Farben \mbox{in colormeta.tex}}}
  \\
\end{tabular}
\end{tcolorbox}

\vspace{0.55em}

% ── 10: Projektstruktur & Build ────────────────────────────────────────────────
\section*{Projektstruktur \& Build-System}

\noindent
\begin{minipage}[t]{0.487\linewidth}

\begin{tcolorbox}[cs, title={Verzeichnisstruktur}]
\ttfamily\small
\begin{tabular}{@{}ll@{}}
  \textbf{config/}       & Design-Tokens, Metadaten \\[2pt]
  \quad colormeta.tex    & Farben, Schriften, Abstände \\
  \quad usermeta.tex     & Autor, Kurs, Matrikel \\
  \quad topicmeta.tex    & Titel, Fach, Beschreibung \\[4pt]
  \textbf{preamble/}     & Technisches Gerüst \\[2pt]
  \quad packages.tex     & Alle \texttt{\textbackslash usepackage} \\
  \quad commands.tex     & Umgebungen \& Befehle \\
  \quad styles.tex       & Layout, Köpfe, TOC \\[4pt]
  \textbf{content/}      & Eigene Inhalte \\[2pt]
  \quad example.tex      & Beispieldatei \\[4pt]
  \textbf{bibliography/} & Literaturquellen \\[2pt]
  \quad sources.bib      & BibTeX-Einträge \\
\end{tabular}
\end{tcolorbox}

\end{minipage}%
\hfill
\begin{minipage}[t]{0.487\linewidth}

\begin{tcolorbox}[cs, title={Makefile-Befehle}]
\begin{tabular}{@{}ll@{}}
  \texttt{make build}     & Einmalig kompilieren \\[4pt]
  \texttt{make watch}     & Automatisch bei Änderungen \\[4pt]
  \texttt{make clean}     & Hilfsdateien löschen \\[4pt]
  \texttt{make distclean} & Alles inkl. PDF löschen \\[4pt]
  \multicolumn{2}{@{}l}{\itshape Alternativ direkt:} \\[2pt]
  \multicolumn{2}{@{}l}{\texttt{latexmk -lualatex main.tex}} \\
\end{tabular}
\end{tcolorbox}

\vspace{0.4em}

\begin{tcolorbox}[cs, title={Überschriften-Hierarchie}]
\begin{tabular}{@{}ll@{}}
  \texttt{\textbackslash section}          & Dunkelblau, Linie \\[4pt]
  \texttt{\textbackslash subsection}       & Abgeschwächtes Blau \\[4pt]
  \texttt{\textbackslash subsubsection}    & Teal/Petrol \\[6pt]
  \multicolumn{2}{@{}l}{\itshape Titelseite generieren:} \\[2pt]
  \multicolumn{2}{@{}l}{\texttt{\textbackslash makelernzetteltitle}} \\
\end{tabular}
\end{tcolorbox}

\end{minipage}

\end{document}
